\label{app:benchmarking_manipuladores}

El análisis comparativo de manipuladores comerciales realizado para establecer valores de referencia en el diseño de los brazos de \emph{Justina} considera nueve modelos representativos divididos en dos categorías: \emph{cobots} de uso industrial y manipuladores de robots de servicio, listados en la Tabla~\ref{tab:bm_brazos_sistemas}.

\Needspace{14\baselineskip}
\begin{longtable}{@{}>{\raggedright\arraybackslash}p{0.35\textwidth} >{\raggedright\arraybackslash}p{0.35\textwidth}@{}}
\caption{Manipuladores considerados en el análisis comparativo.}\label{tab:bm_brazos_sistemas}\\
\toprule
\textbf{Sistema} & \textbf{Marca} \\
\midrule
\endfirsthead
\caption[]{Manipuladores considerados en el análisis comparativo, continuación.}\\
\toprule
\textbf{Sistema} & \textbf{Marca} \\
\midrule
\endhead
\midrule
\multicolumn{2}{r}{\small\emph{Continúa en la siguiente página}}\\
\endfoot
\bottomrule
\endlastfoot
\rowcolor{nprim} \multicolumn{2}{@{}l}{\textbf{\emph{Cobots} de uso industrial}} \\
Panda~\cite{franka_panda} & Franka Emika \\
Gen3~\cite{kinova_gen3} & Kinova \\
UR3e~\cite{ur3e_datasheet} & Universal Robots \\
EBLM~\cite{ti5_eblm} & Ti5 Robot \\
\rowcolor{nprim} \multicolumn{2}{@{}l}{\textbf{Manipuladores de robots de servicio}} \\
Stretch 3~\cite{hello_stretch3} & Hello Robot \\
HSR~\cite{toyota_hsr} & Toyota \\
TIAGo~\cite{pal_tiago} & PAL Robotics \\
G1~\cite{unitree_g1} & Unitree \\
OpenArm~\cite{openarm_hardware} & Enactic \\
\end{longtable}

Las métricas de comparación responden a las necesidades seleccionadas para el diseño de los manipuladores, enumeradas en la Tabla~\ref{tab:bm_brazos_necesidades}.

\Needspace{12\baselineskip}
\begin{longtable}{@{}>{\centering\arraybackslash}p{0.08\textwidth} >{\raggedright\arraybackslash}p{0.7\textwidth}@{}}
\caption{Necesidades seleccionadas para el diseño de los manipuladores.}\label{tab:bm_brazos_necesidades}\\
\toprule
\textbf{\#} & \textbf{Necesidad} \\
\midrule
\endfirsthead
\caption[]{Necesidades seleccionadas para el diseño de los manipuladores, continuación.}\\
\toprule
\textbf{\#} & \textbf{Necesidad} \\
\midrule
\endhead
\midrule
\multicolumn{2}{r}{\small\emph{Continúa en la siguiente página}}\\
\endfoot
\bottomrule
\endlastfoot
1 & El robot manipula artículos domésticos. \\
2 & El robot planea trayectorias del manipulador. \\
3 & El robot toma artículos con la fuerza necesaria. \\
4 & El robot extiende el alcance efectivo de sus manipuladores. \\
5 & El robot manipula equipo para realizar tareas domésticas. \\
6 & El robot es robusto. \\
\end{longtable}

La Tabla~\ref{tab:bm_brazos_metricas} presenta las cinco métricas empleadas en la comparación, junto con sus necesidades asociadas, unidades e importancia en una escala del 1 al 5.

\begin{table}[htbp]
\centering
\caption{Métricas de comparación de los manipuladores.}
\label{tab:bm_brazos_metricas}
\renewcommand{\arraystretch}{1.3}
\begin{tabularx}{\textwidth}{cc>{\centering\arraybackslash}Xcc}
\hline
\textbf{\# Métrica} & \textbf{\# Necesidad} & \textbf{Métrica} & \textbf{Importancia} & \textbf{Unidad} \\
\hline
1 & 1, 5, 6 & Capacidad de carga.     & 5 & {[}kg{]} \\
3 & 2, 4, 5 & Alcance.                & 5 & {[}mm{]} \\
5 & 4, 5    & Grados de libertad.     & 4 & {[}1{]}  \\
6 & 2, 6    & Repetibilidad.          & 4 & {[}mm{]} \\
7 & 5, 6    & Grado de protección IP. & 3 & {[}1{]}  \\
\hline
\end{tabularx}
\end{table}

La Tabla~\ref{tab:bm_brazos_especificaciones} reúne los valores de estas métricas extraídos de las hojas de especificaciones disponibles para cada sistema.

\Needspace{18\baselineskip}
\begingroup\small\setlength{\tabcolsep}{4pt}
\begin{longtable}{@{}>{\raggedright\arraybackslash}p{0.15\textwidth} >{\centering\arraybackslash}p{0.14\textwidth} >{\centering\arraybackslash}p{0.14\textwidth} >{\centering\arraybackslash}p{0.14\textwidth} >{\centering\arraybackslash}p{0.14\textwidth} >{\centering\arraybackslash}p{0.14\textwidth}@{}}
\caption{Especificaciones de los manipuladores comerciales.}\label{tab:bm_brazos_especificaciones}\\
\toprule
\textbf{Sistema} & \textbf{Capacidad de carga [kg]} & \textbf{Alcance [mm]} & \textbf{Grados de libertad [1]} & \textbf{Repeti\-bilidad [mm]} & \textbf{Grado de protección IP [1]} \\
\midrule
\endfirsthead
\caption[]{Especificaciones de los manipuladores comerciales, continuación.}\\
\toprule
\textbf{Sistema} & \textbf{Capacidad de carga [kg]} & \textbf{Alcance [mm]} & \textbf{Grados de libertad [1]} & \textbf{Repeti\-bilidad [mm]} & \textbf{Grado de protección IP [1]} \\
\midrule
\endhead
\midrule
\multicolumn{6}{r}{\small\emph{Continúa en la siguiente página}}\\
\endfoot
\bottomrule
\endlastfoot
\rowcolor{nprim} \multicolumn{6}{@{}l}{\textbf{\emph{Cobots} de uso industrial}} \\
Panda & 3 & 855 & 7 & $\pm$0.1 & IP30 \\
Gen3 & 2 & 902 & 7 & $\pm$1 & IP33 \\
UR3e & 3 & 500 & 6 & $\pm$0.03 & IP54 \\
EBLM & 3 & 600 & 6 & $\pm$0.05 & IP50 \\
\rowcolor{nprim} \multicolumn{6}{@{}l}{\textbf{Manipuladores de robots de servicio}} \\
Stretch 3 & 3 & 510 & 4 & --- & --- \\
HSR & 1.2 & 600 & 4 & --- & --- \\
TIAGo & 2.8 & 870 & 7 & --- & --- \\
G1 & 2 & 450 & 5 & --- & --- \\
OpenArm & 4.1 & 633 & 7 & --- & --- \\
\end{longtable}
\endgroup

Con base en estos valores se elaboró una matriz de comparación en la que se evaluaron los nueve sistemas considerados. Las métricas se agruparon según las necesidades a las que responden: planeación de trayectorias, alcance efectivo, tareas domésticas y robustez. Para cada grupo, la Tabla~\ref{tab:bm_brazos_qfd} presenta la importancia $w_j$ de cada métrica $j$ definida en la Tabla~\ref{tab:bm_brazos_metricas}, la calificación normalizada $c_{ij} \in [0,1]$ de cada sistema $i$, donde 1 corresponde al mejor valor del conjunto evaluado, y el puntaje total $T_i$ del grupo. Este puntaje se expresa en una escala de 0 a 5, redondeado al entero más cercano, mediante la relación:

\[
T_i = 5 \, \frac{\sum_j w_j \, c_{ij}}{\sum_j w_j}.
\]

\Needspace{20\baselineskip}
\begingroup\scriptsize\setlength{\tabcolsep}{1pt}
\begin{longtable}{@{}>{\raggedright\arraybackslash\leftskip=1.2em\hyphenpenalty=10000\exhyphenpenalty=10000}p{0.15\textwidth} >{\centering\arraybackslash}p{0.05\textwidth} >{\centering\arraybackslash}p{0.078\textwidth} >{\centering\arraybackslash}p{0.078\textwidth} >{\centering\arraybackslash}p{0.078\textwidth} >{\centering\arraybackslash}p{0.078\textwidth} >{\centering\arraybackslash}p{0.078\textwidth} >{\centering\arraybackslash}p{0.078\textwidth} >{\centering\arraybackslash}p{0.078\textwidth} >{\centering\arraybackslash}p{0.078\textwidth} >{\centering\arraybackslash}p{0.10\textwidth}@{}}
\caption{Matriz de comparación de los manipuladores comerciales.}\label{tab:bm_brazos_qfd}\\
\toprule
{\leftskip=0pt\scriptsize\textbf{Métrica}} & {\scriptsize\textbf{Imp.}} & {\scriptsize\textbf{Panda}} & {\scriptsize\textbf{Gen3}} & {\scriptsize\textbf{UR3e}} & {\scriptsize\textbf{EBLM}} & {\scriptsize\textbf{Stretch 3}} & {\scriptsize\textbf{HSR}} & {\scriptsize\textbf{TIAGo}} & {\scriptsize\textbf{G1}} & {\scriptsize\textbf{OpenArm}} \\
\midrule
\endfirsthead
\caption[]{Matriz de comparación de los manipuladores comerciales, continuación.}\\
\toprule
{\leftskip=0pt\scriptsize\textbf{Métrica}} & {\scriptsize\textbf{Imp.}} & {\scriptsize\textbf{Panda}} & {\scriptsize\textbf{Gen3}} & {\scriptsize\textbf{UR3e}} & {\scriptsize\textbf{EBLM}} & {\scriptsize\textbf{Stretch 3}} & {\scriptsize\textbf{HSR}} & {\scriptsize\textbf{TIAGo}} & {\scriptsize\textbf{G1}} & {\scriptsize\textbf{OpenArm}} \\
\midrule
\endhead
\midrule
\multicolumn{11}{r}{\small\emph{Continúa en la siguiente página}}\\
\endfoot
\bottomrule
\endlastfoot
\rowcolor{nprim} \multicolumn{11}{@{}l}{\textbf{Planeación de trayectorias}} \\*
Alcance & 5 & 0.8 & 1.0 & 0.2 & 0.6 & 0.4 & 0.6 & 0.8 & 0.0 & 0.7 \\*
Grados de libertad & 4 & 1.0 & 1.0 & 0.8 & 0.8 & 0.4 & 0.4 & 1.0 & 0.6 & 1.0 \\*
Repetibilidad & 4 & 0.7 & 0.6 & 1.0 & 0.8 & 0.0 & 0.0 & 0.0 & 0.0 & 0.0 \\*
\rowcolor{nsec} \textbf{Total} & & \textbf{4} & \textbf{4} & \textbf{3} & \textbf{4} & \textbf{1} & \textbf{2} & \textbf{3} & \textbf{1} & \textbf{3} \\
\rowcolor{nprim} \multicolumn{11}{@{}l}{\textbf{Alcance efectivo}} \\*
Alcance & 5 & 0.8 & 1.0 & 0.2 & 0.6 & 0.4 & 0.6 & 0.8 & 0.0 & 0.7 \\*
Grados de libertad & 4 & 1.0 & 1.0 & 0.8 & 0.8 & 0.4 & 0.4 & 1.0 & 0.6 & 1.0 \\*
\rowcolor{nsec} \textbf{Total} & & \textbf{4} & \textbf{5} & \textbf{2} & \textbf{3} & \textbf{2} & \textbf{3} & \textbf{4} & \textbf{1} & \textbf{4} \\
\rowcolor{nprim} \multicolumn{11}{@{}l}{\textbf{Tareas domésticas}} \\*
Capacidad de carga & 5 & 0.8 & 0.6 & 0.8 & 0.8 & 0.8 & 0.4 & 0.8 & 0.6 & 1.0 \\*
Alcance & 5 & 0.8 & 1.0 & 0.2 & 0.6 & 0.4 & 0.6 & 0.8 & 0.0 & 0.7 \\*
Grados de libertad & 4 & 1.0 & 1.0 & 0.8 & 0.8 & 0.4 & 0.4 & 1.0 & 0.6 & 1.0 \\*
Repetibilidad & 4 & 0.7 & 0.6 & 1.0 & 0.8 & 0.0 & 0.0 & 0.0 & 0.0 & 0.0 \\*
Grado de protección IP & 3 & 0.4 & 0.6 & 1.0 & 0.8 & 0.0 & 0.0 & 0.0 & 0.0 & 0.0 \\*
\rowcolor{nsec} \textbf{Total} & & \textbf{4} & \textbf{4} & \textbf{4} & \textbf{4} & \textbf{2} & \textbf{2} & \textbf{3} & \textbf{1} & \textbf{3} \\
\rowcolor{nprim} \multicolumn{11}{@{}l}{\textbf{Robustez}} \\*
Capacidad de carga & 5 & 0.8 & 0.6 & 0.8 & 0.8 & 0.8 & 0.4 & 0.8 & 0.6 & 1.0 \\*
Repetibilidad & 4 & 0.7 & 0.6 & 1.0 & 0.8 & 0.0 & 0.0 & 0.0 & 0.0 & 0.0 \\*
Grado de protección IP & 3 & 0.4 & 0.6 & 1.0 & 0.8 & 0.0 & 0.0 & 0.0 & 0.0 & 0.0 \\*
\rowcolor{nsec} \textbf{Total} & & \textbf{3} & \textbf{3} & \textbf{5} & \textbf{4} & \textbf{2} & \textbf{1} & \textbf{2} & \textbf{1} & \textbf{2} \\
\end{longtable}
\endgroup

Los \emph{cobots} de uso industrial combinan alcance, grados de libertad y precisión. Todos poseen seis o siete grados de libertad, alcances de 500 a 902~mm y capacidad de carga de 2 a 3~kg, y son los únicos en reportar repetibilidad, de $\pm$0.03 a $\pm$1~mm, y grado de protección IP, de IP30 a IP54. Sobresalen el Gen3 con 902~mm de alcance y el UR3e con $\pm$0.03~mm de repetibilidad y grado de protección IP54.

Los manipuladores de robots de servicio omiten datos de repetibilidad y grado de protección IP. Mientras el Stretch~3, el HSR y el G1 se limitan a alcances de 450 a 600~mm y cuatro o cinco grados de libertad, el TIAGo y el OpenArm igualan la movilidad de los \emph{cobots} con siete grados de libertad y alcances de 870 y 633~mm, destacando el OpenArm con la mayor capacidad de carga, de 4.1~kg.

Ningún sistema domina todos los grupos. El Gen3 obtiene 5 puntos en alcance efectivo y 4 en planeación de trayectorias y tareas domésticas, pero baja a 3 en robustez. El Panda y el EBLM logran el perfil más equilibrado con 4 puntos en tres grupos; el UR3e lidera en robustez con 5 y obtiene 4 en tareas domésticas, aunque solo 2 en alcance efectivo. Entre los manipuladores de servicio destacan el TIAGo y el OpenArm, ambos con 4 en alcance efectivo y 3 en planeación de trayectorias y tareas domésticas.
